\section{System Overview}
\label{sec:system}

SHI is architected around a dual-stage interactive retrieval loop that connects
semantic narrative structuring with occurrence-level temporal steering. To
support multi-event search over large-scale video collections, offline visual
frames are pre-indexed using SigLIP2~\cite{tschannen2025siglip2}, contextual
and speech metadata are encoded with BGE-M3~\cite{chen2024bgem3}, and BM25
~\cite{robertson2009bm25} provides lexical scoring. Building upon these base
representations, the primary scientific contribution of SHI lies in how its
pre-retrieval and post-retrieval interaction modules interoperate through a
shared state representation, as shown in the interface overview in
Fig.~\ref{fig:system-ui} and the retrieval pipeline in Fig.~\ref{fig:layer2}.

\begin{figure}[!t]
  \centering
  \includegraphics[width=0.78\textwidth]{001_system_ui.png}
  \caption{SHI interface for multi-event KIS. Candidate videos and their
  decoded event paths are shown on the left, while the Query Hypothesis editor
  on the right allows the event structure to be inspected and revised before
  retrieval.}
  \label{fig:system-ui}
\end{figure}

\subsection{Interaction Loop and Module Interplay}
\label{sec:interaction-interplay}

The two interaction stages in SHI form a closed human-in-the-loop workflow.
\emph{First}, in the pre-retrieval phase, the Query Hypothesis converts the
unstructured narrative into an explicit, ordered sequence of searchable events;
rather than executing this decomposition blindly, the user verifies event
boundaries and temporal sequencing, binding visual reference images where
applicable to ensure multimodal scoring is driven by an inspected semantic
hypothesis.
\emph{Second}, the retrieval pipeline acts as the decoding bridge: each accepted
event queries the multimodal indexes independently, and component similarities
are calibrated and combined via event-adaptive fusion, yielding an event--frame
score matrix $S_v \in \mathbb{R}^{N \times T_v}$ for each candidate video $v$. A
dynamic-programming decoder searches this matrix for a chronologically ordered
path assigning one frame occurrence per event under temporal gap penalties,
producing an initial ranking of candidate videos.
\emph{Third}, in the post-retrieval phase, EventTrail exposes the decoded alignment of
a selected candidate as a coherent temporal path rather than disjoint frames;
when an event occurrence is misidentified, the user applies localized actions
(keep, use, or reject), injecting anchor or exclusion constraints directly into
the path decoder.
\emph{Finally}, the system supports cross-stage recovery: if candidate inspection
reveals that the intended target cannot be matched under the current event
formulation, the operator returns to the Query Hypothesis to restructure the
underlying narrative before executing a revised corpus-level search.

\subsection{Evidence Snapshot as the Decoupling Mechanism}
\label{sec:evidence-snapshots}

Interactive steering in live competition demands sub-second feedback. SHI
achieves this by freezing the complete retrieval state of top-ranked videos
into session-scoped Evidence Snapshots. Each snapshot preserves the calibrated
event--frame score matrix $S_v$, decoded candidate paths, event order, and
decoder hyperparameters. Consequently, occurrence-level corrections in EventTrail
operate entirely on retained evidence via constrained dynamic programming. This
effectively decouples local temporal path re-alignment ($O(N \cdot T_v)$) from
expensive global corpus retrieval ($O(|\mathcal{V}|)$), enabling instantaneous
feedback. Once the correct target interval is identified, the submission is
dispatched directly through the integrated DRES client~\cite{rossetto2021dres}.